\section{Reproducibility Details}

The code package contains a lightweight NumPy implementation. The primary entry point is:
\begin{verbatim}
python -m qprotohop
\end{verbatim}
The runner supports synthetic datasets, local \texttt{.npz} datasets with \texttt{x/y}, \texttt{x\_train/y\_train/x\_test/y\_test}, or \texttt{xtr/ytr/xte/yte} keys, and optional torchvision exports for MNIST and Fashion-MNIST.

\section{Real-Data Commands}

The real-data experiments in this manuscript were generated with:
\begin{verbatim}
scripts/run_real_cifar4_final.ps1
scripts/run_real_cifar4_shot_sweep.ps1
scripts/run_real_cifar4_comm_budget.ps1
scripts/run_real_mnist_fashion_final.ps1
scripts/run_real_mnist_fashion_hop_ablation.ps1
\end{verbatim}
Tables and figures are produced by:
\begin{verbatim}
scripts/report_real_datasets.py
scripts/report_real_cifar4.py
scripts/plot_real_figures.py
scripts/export_real_latex_tables.py
\end{verbatim}

\section{Synthetic Commands}

The synthetic schema, high-order, and communication experiments are generated by:
\begin{verbatim}
scripts/run_schema_triptych.ps1
scripts/run_hop_ablation.ps1
scripts/run_iso_communication.ps1
scripts/run_comm_budget_curve.ps1
scripts/run_noise_heterogeneity_sweep.ps1
scripts/run_schema_mismatch_sweep.ps1
\end{verbatim}

\section{Quantum-Backend and Strong-Baseline Commands}

The PennyLane/Qiskit, CHOP, and stronger-baseline checks are generated by:
\begin{verbatim}
scripts/run_qproto_masked_followups.ps1
scripts/run_qproto_masked_hop_followups.ps1
scripts/run_qproto_chop_followups.ps1
scripts/run_chop_k_sweeps.ps1
scripts/run_chop_key_policy_controls.ps1
scripts/run_strong_baselines.ps1
scripts/run_schema_head_baselines.ps1
scripts/run_drift_baselines.ps1
scripts/run_quantum_param_vqc_baseline.ps1
scripts/run_if_fusion_controls.ps1
scripts/run_if_scientific_fusion.ps1
scripts/run_scale_overlap_sweeps.ps1
scripts/run_hardware_ready_demo.ps1
scripts/run_qiskit_aer_noise_sweep.ps1
\end{verbatim}

The paired significance/effect-size table is generated from the shared summary file with:
\begin{verbatim}
scripts/report_if_significance.py
\end{verbatim}

The real non-quantum IF benchmarks are generated by:
\begin{verbatim}
scripts/download_if_real_datasets.ps1
scripts/download_mhealth_dataset.ps1
scripts/download_hydraulic_dataset.ps1
scripts/download_pamap2_dataset.ps1
scripts/run_uci_har_if_benchmark.ps1
scripts/run_mfeat_if_benchmark.ps1
scripts/run_mhealth_if_benchmark.ps1
scripts/run_wdbc_if_benchmark.ps1
scripts/run_gas_drift_if_benchmark.ps1
scripts/run_hydraulic_if_benchmark.ps1
scripts/run_pamap2_if_benchmark.ps1
scripts/screen_gas_chop_configs.py
scripts/export_adaptive_chop_ablation.py
scripts/report_adaptive_chop_gate_choices.py
\end{verbatim}

\section{Implementation Notes}

The HOP branch is applied to schema-adapted readout values rather than to the value-and-mask vector. This avoids turning missingness masks into high-order class features. The low-order branch still receives both values and masks, so missing observables remain distinguishable from zero-valued expectations.

The reported communication is computed as the number of transmitted scalar statistics times four bytes. The count includes class prototypes, class counts, high-order sketches when used, and small metadata fields. Public schema definitions and anchors are treated as protocol setup objects rather than per-round payloads.

\section{Additional Experiment Tables}

This appendix keeps QFL-specific and diagnostic results out of the IF main narrative while preserving the evidence needed for reproducibility and specialist review.

\subsection{Claim-to-Evidence Map}

\begin{table*}[t]
\centering
\caption{How the supplementary evidence supports the main scoped claims. The table is intended to make clear which results are headline evidence and which results are boundary checks.}
\label{tab:claim-evidence-map}
\scriptsize
\resizebox{\textwidth}{!}{%
\begin{tabular}{p{0.26\textwidth}p{0.36\textwidth}p{0.30\textwidth}}
\toprule
Scoped claim & Primary evidence & Boundary or caveat \\
\midrule
Semantic comparability is necessary before aggregation & No-schema, forced-canonical, wrong-schema, zero-fill, mask-only, shared-view, and schema-correct rows in the main image and IF-control tables & Shared-view can be strong when the common intersection already carries the signal; this is treated as a boundary case, not a failure of the protocol view \\
Compact value/count fusion is useful under incomplete but identified sources & UCI HAR, Multiple Features, MHEALTH, PAMAP2, Hydraulic Systems, and WDBC equal-byte tables; paired tests in Table~\ref{tab:if-significance-full} & Full-dimensional HeMIS, MLP, imputation, and late-fusion rows are stronger boundary references when centralized feature-level training is allowed \\
HOP/CHOP is conditional, not universal & Covariance-sensing, MHEALTH, PennyLane high-order readout, CHOP sweeps, and key-policy diagnostics & HAR, MFeat, PAMAP2, WDBC, and Gas Drift show that low-order CProto or adaptive fallback is safer when the signal is low-order or imputation-friendly \\
Selected keys should be tested against random compression & Random-key CProto/HOP controls in MFeat, MHEALTH, WDBC, and CHOP key-policy tables & MFeat random-key low-order CProto is strong; that dataset supports schema-aware low-order fusion rather than selected-key superiority \\
Quantum-readout compatibility is engineering evidence, not hardware validation & Qiskit Aer count-derived robustness sweep, PennyLane high-order readouts, and hardware-ready path table & No paid IBM hardware data are used; hardware validation remains future work \\
Communication fairness is explicit & Strict equal-byte tables, payload accounting notes, and protocol manifest & Public schema definitions and anchors are treated as setup objects rather than per-round payloads \\
\bottomrule
\end{tabular}%
}
\end{table*}

\subsection{Protocol Manifest}

\begin{table*}[t]
\centering
\caption{Protocol manifest for the main experiments. Anchors are public/unlabeled and are used only for preprocessing, imputation, or key selection; test labels are never used for schema construction or key selection.}
\label{tab:protocol-manifest}
\scriptsize
\resizebox{\textwidth}{!}{%
\begin{tabular}{lcccccc}
\toprule
Experiment & Seeds & Clients & Source dictionary & Client schema sampling & Public anchors & Tuned hyperparameters \\
\midrule
Full-class images & 5 & 20 & simulated observable dictionary from image features & heterogeneous observable subsets/order & train-only public subset & \(K,h,\lambda\) fixed by preliminary grid \\
UCI HAR & 5 & 20 & feature-name sensor groups & 1 common group + random private groups & train-only anchor features & PCA rank, \(K,h,\lambda\), torch epochs fixed before reporting \\
UCI Multiple Features & 10 & 20 & six feature families & 1 common weak view + 2 private views & train-only anchor features & source-balanced score, \(K,h,\lambda\), torch epochs fixed before reporting \\
MHEALTH & 10 & 20 & chest/ECG/ankle/arm wearable sources & subject-held-out windows, 1 common source + private sources & train-only anchor windows & \(K,h,\lambda\), window size, and torch epochs fixed before reporting \\
PAMAP2 & 5 & 20 & heart-rate/hand/chest/ankle wearable sources & subject-held-out windows, 1 common source + private sources & train-only anchor windows & \(K,h,\lambda\), adaptive margin, window size, and torch epochs fixed before reporting \\
Hydraulic Systems & 5 & 20 & named industrial sensor blocks & cooler-state clients with shared/private sensor groups & train-only anchor records & \(K,h,\lambda\), HeMIS/MLP epochs, and equal-byte payload fixed before reporting \\
WDBC & 10 & 20 & ten biomedical measurement families & 1 common family + random private families & train-only anchor features & same non-quantum grid as HAR/MFeat \\
Gas Drift boundary check & 5 & 20 & gas-sensor blocks & random and temporal source splits & train-only anchor records & cross-sensor HOP screen and low-order baselines fixed before reporting \\
Covariance sensing & 5 & 20 & controlled source dictionary & heterogeneous low-overlap schemas & generated public anchors & \(K,h,\lambda\) selected by communication budget \\
Qiskit/PennyLane checks & 3--5 & 20 & circuit-generated observable readouts & heterogeneous basis/subset/depth settings & circuit-generated public anchors & no test-label tuning; Aer/PennyLane settings fixed \\
\bottomrule
\end{tabular}%
}
\end{table*}


The seed counts in Table~\ref{tab:protocol-manifest} are the counts used for the reported means and paired tests. Strict equal-byte tables below use the same seed identifiers, payload formulas, and public-anchor splits as the corresponding main-text rows. Random-key controls are regenerated from the experiment seed only and do not use validation or test statistics for key selection.

\subsection{Four-Class Image Pilots}

\begin{table*}[t]
\centering
\caption{Real image datasets under heterogeneous quantum readout. Results are means with 95\% confidence intervals over 5 seeds.}
\label{tab:real-image-main}
\begin{tabular}{llccc}
\toprule
Dataset & Method & Acc. (95\% CI) & Bal. acc. & Bytes \\
\midrule
\multirow{8}{*}{MNIST-4}
& QProto-HOP & 0.709 $\pm$ 0.054 & 0.707 & 4628 \\
& QProto prototype & 0.709 $\pm$ 0.054 & 0.707 & 3092 \\
& No schema & 0.441 $\pm$ 0.029 & 0.437 & 3092 \\
& Forced canonical & 0.437 $\pm$ 0.039 & 0.434 & 3092 \\
& Wrong schema & 0.280 $\pm$ 0.023 & 0.278 & 4628 \\
& FedAvg forced & 0.374 $\pm$ 0.078 & 0.368 & 3088 \\
& FedProx forced & 0.374 $\pm$ 0.079 & 0.368 & 3088 \\
& Classical kernel & 0.671 $\pm$ 0.026 & 0.667 & 3092 \\
\midrule
\multirow{8}{*}{Fashion-4}
& QProto-HOP & 0.746 $\pm$ 0.037 & 0.746 & 4628 \\
& QProto prototype & 0.745 $\pm$ 0.038 & 0.745 & 3092 \\
& No schema & 0.431 $\pm$ 0.024 & 0.430 & 3092 \\
& Forced canonical & 0.429 $\pm$ 0.022 & 0.428 & 3092 \\
& Wrong schema & 0.299 $\pm$ 0.027 & 0.298 & 4628 \\
& FedAvg forced & 0.336 $\pm$ 0.064 & 0.339 & 3088 \\
& FedProx forced & 0.335 $\pm$ 0.064 & 0.338 & 3088 \\
& Classical kernel & 0.663 $\pm$ 0.017 & 0.663 & 3092 \\
\midrule
\multirow{8}{*}{CIFAR-4}
& QProto-HOP & 0.486 $\pm$ 0.052 & 0.485 & 4628 \\
& QProto prototype & 0.486 $\pm$ 0.052 & 0.485 & 3092 \\
& No schema & 0.339 $\pm$ 0.017 & 0.339 & 3092 \\
& Forced canonical & 0.339 $\pm$ 0.019 & 0.339 & 3092 \\
& Wrong schema & 0.268 $\pm$ 0.015 & 0.267 & 4628 \\
& FedAvg forced & 0.298 $\pm$ 0.046 & 0.299 & 3088 \\
& FedProx forced & 0.298 $\pm$ 0.046 & 0.299 & 3088 \\
& Classical kernel & 0.411 $\pm$ 0.010 & 0.411 & 3092 \\
\bottomrule
\end{tabular}
\end{table*}


\begin{figure}[t]
    \centering
    \includegraphics[width=\linewidth]{real_image_dataset_accuracy.pdf}
    \caption{Four-class image diagnostics under heterogeneous quantum readout. These results motivated the full-class experiments in the main text.}
    \label{fig:real-image}
\end{figure}

\subsection{FL and QFL Baselines}

\begin{table*}[t]
\centering
\caption{Strong FL baselines under the same heterogeneous-readout protocols. Full-class rows use five seeds; Qiskit Aer and PennyLane rows use three seeds. Bytes are per client per round.}
\label{tab:strong-baselines}
\scriptsize
\renewcommand{\arraystretch}{0.93}
\begin{tabular}{llccc}
\toprule
Setting & Method & Acc. & 95\% CI & Bytes \\
\midrule
MNIST-10 & FedAdam forced & 0.229 & 0.011 & 7720 \\
MNIST-10 & FedAdam schema & 0.456 & 0.033 & 7720 \\
MNIST-10 & FedProto schema & 0.412 & 0.029 & 7724 \\
MNIST-10 & CProto, \(K=96\) & 0.738 & 0.030 & 7724 \\
MNIST-10 & CHOP, \(K=96\) & 0.739 & 0.030 & 11564 \\
MNIST-10 & \method{}-Masked-HOP & 0.744 & 0.027 & 24364 \\
\midrule
Fashion-10 & FedAdam forced & 0.254 & 0.019 & 7720 \\
Fashion-10 & FedAdam schema & 0.544 & 0.033 & 7720 \\
Fashion-10 & FedProto schema & 0.522 & 0.037 & 7724 \\
Fashion-10 & CProto, \(K=96\) & 0.660 & 0.027 & 7724 \\
Fashion-10 & CHOP, \(K=96\) & 0.660 & 0.027 & 11564 \\
Fashion-10 & \method{}-Masked-HOP & 0.662 & 0.026 & 24364 \\
\midrule
CIFAR-10 & FedAdam forced & 0.166 & 0.012 & 7720 \\
CIFAR-10 & FedAdam schema & 0.281 & 0.017 & 7720 \\
CIFAR-10 & FedProto schema & 0.268 & 0.019 & 7724 \\
CIFAR-10 & CProto, \(K=96\) & 0.450 & 0.005 & 7724 \\
CIFAR-10 & CHOP, \(K=96\) & 0.451 & 0.004 & 11564 \\
CIFAR-10 & \method{}-Masked-HOP & 0.454 & 0.004 & 24364 \\
\midrule
Qiskit Aer MNIST-4 & FedAdam forced & 0.379 & 0.034 & 528 \\
Qiskit Aer MNIST-4 & FedAdam schema & 0.426 & 0.178 & 528 \\
Qiskit Aer MNIST-4 & FedProto schema & 0.428 & 0.197 & 532 \\
Qiskit Aer MNIST-4 & CProto, \(K=8\) & 0.576 & 0.045 & 276 \\
Qiskit Aer MNIST-4 & CHOP, \(K=8\) & 0.579 & 0.043 & 532 \\
Qiskit Aer MNIST-4 & \method{}-Masked-HOP & 0.597 & 0.065 & 660 \\
\midrule
PennyLane high-order & FedAdam forced & 0.359 & 0.037 & 1552 \\
PennyLane high-order & FedAdam schema & 0.317 & 0.021 & 1552 \\
PennyLane high-order & FedProto schema & 0.285 & 0.029 & 1556 \\
PennyLane high-order & CProto, \(K=24\) & 0.403 & 0.061 & 788 \\
PennyLane high-order & CHOP, \(K=24\) & 0.602 & 0.060 & 2324 \\
PennyLane high-order & \method{}-Masked-HOP & 0.550 & 0.009 & 3476 \\
\bottomrule
\end{tabular}
\end{table*}

\begin{table*}[t]
\centering
\caption{Drift-control FL baselines on quantum-backend readout tasks. Values are mean accuracy over three seeds; bytes are per client per round. SCAFFOLD transmits an additional control-variate update, so its communication is higher than ordinary head aggregation.}
\label{tab:drift-baselines}
\begin{tabular}{llccc}
\toprule
Setting & Method & Acc. & 95\% CI & Bytes \\
\midrule
Qiskit Aer MNIST-4 & FedAdam schema & 0.426 & 0.178 & 528 \\
Qiskit Aer MNIST-4 & FedProto schema & 0.428 & 0.197 & 532 \\
Qiskit Aer MNIST-4 & SCAFFOLD schema & 0.357 & 0.191 & 1056 \\
Qiskit Aer MNIST-4 & FedDyn schema & 0.418 & 0.188 & 528 \\
Qiskit Aer MNIST-4 & CProto & 0.576 & 0.045 & 276 \\
Qiskit Aer MNIST-4 & CHOP & 0.579 & 0.043 & 532 \\
Qiskit Aer MNIST-4 & \method{}-Masked-HOP & \textbf{0.597} & 0.065 & 660 \\
\midrule
PennyLane high-order & FedAdam schema & 0.316 & 0.021 & 1552 \\
PennyLane high-order & FedProto schema & 0.285 & 0.029 & 1556 \\
PennyLane high-order & SCAFFOLD schema & 0.283 & 0.009 & 3104 \\
PennyLane high-order & FedDyn schema & 0.299 & 0.011 & 1552 \\
PennyLane high-order & CProto & 0.403 & 0.061 & 788 \\
PennyLane high-order & CHOP & \textbf{0.602} & 0.061 & 2324 \\
PennyLane high-order & \method{}-Masked-HOP & 0.550 & 0.009 & 3476 \\
\bottomrule
\end{tabular}
\end{table*}

\begin{table}[t]
\centering
\caption{Trainable PennyLane VQC parameter-aggregation baseline. This baseline uses a shared ansatz and canonical readout, so it tests standard quantum-parameter FL rather than heterogeneous-readout prototype aggregation.}
\label{tab:quantum-param-vqc}
\begin{tabular}{lrrrr}
\toprule
Method & Acc. & 95\% CI & Bal. acc. & Bytes \\
\midrule
VQC init & 0.243 & 0.085 & 0.235 & 128 \\
VQC FedAvg & 0.535 & 0.084 & 0.514 & 128 \\
VQC FedProx & 0.534 & 0.082 & 0.513 & 128 \\
VQC centralized & 0.520 & 0.063 & 0.500 & 128 \\
\bottomrule
\end{tabular}
\end{table}


These tables compare the protocol against FedAvg/FedProx/FedAdam, SCAFFOLD/FedDyn, FedProto, and trainable PennyLane VQC parameter aggregation. They are important specialist baselines, but they address QFL optimization rather than the central IF claim, so they are kept in the appendix.

\subsection{Scale, Overlap, and Backend Pilots}

\begin{table*}[t]
\centering
\caption{Scale and observable-overlap stress tests. Values are mean accuracy over three seeds. Detailed confidence intervals are reported in the released experiment logs.}
\label{tab:scale-overlap}
\resizebox{\textwidth}{!}{%
\begin{tabular}{llcccccc}
\toprule
Setting & Stress & FedAdam & FedProto & Shared & CProto & CHOP & Masked-HOP \\
\midrule
MNIST-10 & 50 clients & 0.518 & 0.532 & 0.702 & 0.745 & 0.747 & 0.752 \\
MNIST-10 & 100 clients & 0.553 & 0.584 & 0.701 & 0.746 & 0.748 & 0.753 \\
Fashion-10 & 50 clients & 0.590 & 0.602 & 0.660 & 0.669 & 0.669 & 0.672 \\
Fashion-10 & 100 clients & 0.614 & 0.620 & 0.660 & 0.670 & 0.671 & 0.673 \\
CIFAR-10 & 50 clients & 0.320 & 0.325 & 0.424 & 0.449 & 0.449 & 0.453 \\
CIFAR-10 & 100 clients & 0.342 & 0.348 & 0.424 & 0.450 & 0.451 & 0.455 \\
\midrule
MNIST-10 & 20\% overlap & 0.173 & 0.158 & 0.555 & 0.688 & 0.689 & 0.737 \\
MNIST-10 & 100\% overlap & 0.681 & 0.719 & 0.718 & 0.744 & 0.745 & 0.745 \\
Fashion-10 & 20\% overlap & 0.238 & 0.220 & 0.598 & 0.651 & 0.651 & 0.670 \\
Fashion-10 & 100\% overlap & 0.673 & 0.668 & 0.667 & 0.670 & 0.672 & 0.671 \\
CIFAR-10 & 20\% overlap & 0.148 & 0.143 & 0.325 & 0.424 & 0.425 & 0.453 \\
CIFAR-10 & 100\% overlap & 0.420 & 0.438 & 0.438 & 0.459 & 0.460 & 0.460 \\
\midrule
PennyLane high-order & 20\% overlap & 0.299 & 0.280 & 0.425 & 0.385 & 0.593 & 0.558 \\
PennyLane high-order & 50\% overlap & 0.314 & 0.291 & 0.475 & 0.433 & 0.662 & 0.616 \\
PennyLane high-order & 80\% overlap & 0.429 & 0.394 & 0.597 & 0.440 & 0.691 & 0.662 \\
\bottomrule
\end{tabular}%
}
\end{table*}

\begin{table*}[t]
\centering
\caption{Quantum-backend compatibility checks after adding coverage-aware aggregation. PennyLane uses precomputed VQC readout features; Qiskit Aer uses noisy finite-shot simulator readouts.}
\label{tab:quantum-backend-masked}
\resizebox{\textwidth}{!}{%
\begin{tabular}{lcccc}
\toprule
Setting & Method & Acc. & 95\% CI & Bytes \\
\midrule
PennyLane MNIST-4 & \method{}-Masked & 0.733 & 0.063 & 1940 \\
PennyLane MNIST-4 & \method{}-Masked-HOP & 0.734 & 0.061 & 2708 \\
PennyLane MNIST-4 & low-order prototype & 0.413 & -- & 1556 \\
PennyLane private-signal & \method{}-Masked & 0.748 & 0.032 & 1940 \\
PennyLane private-signal & \method{}-Masked-HOP & 0.750 & 0.029 & 2708 \\
PennyLane private-signal & no-schema & 0.405 & 0.041 & 1556 \\
PennyLane private-signal & shared-observable & 0.357 & 0.013 & 1556 \\
\midrule
Qiskit Aer MNIST-4 & \method{}-Masked & 0.593 & 0.060 & 404 \\
Qiskit Aer MNIST-4 & \method{}-Masked-HOP & 0.597 & 0.065 & 660 \\
Qiskit Aer MNIST-4 & CHOP, \(K=8\) & 0.579 & 0.043 & 532 \\
Qiskit Aer MNIST-4 & CProto, \(K=8\) & 0.576 & 0.046 & 276 \\
Qiskit Aer MNIST-4 & shared-observable & 0.558 & 0.083 & 532 \\
Qiskit Aer MNIST-4 & low-order prototype & 0.427 & 0.238 & 532 \\
Qiskit Aer MNIST-4 & no-schema & 0.375 & 0.155 & 532 \\
\bottomrule
\end{tabular}%
}
\end{table*}

\begin{table*}[t]
\centering
\scriptsize
\caption{Qiskit Aer robustness sweep with count-derived readout expectations. Values are mean accuracy over five seeds. The QProto runner adds no extra synthetic readout noise here; finite shots and noise are already present in the Aer-generated measurement counts.}
\label{tab:qiskit-noise-sweep}
\resizebox{\textwidth}{!}{%
\begin{tabular}{lcccccccc}
\toprule
Setting & Wrong & No schema & FedAdam & FedProto & Shared & CProto & CHOP & Masked-HOP \\
\midrule
Clean, depth 1, 1024 shots & 0.320 & 0.360 & 0.438 & 0.454 & 0.518 & 0.568 & 0.572 & \textbf{0.580} \\
Low shot, 128 shots & 0.309 & 0.340 & 0.398 & 0.410 & 0.466 & 0.546 & \textbf{0.550} & 0.547 \\
High readout error & 0.314 & 0.356 & 0.425 & 0.437 & 0.496 & 0.573 & \textbf{0.583} & 0.572 \\
High depolarizing noise & 0.312 & 0.352 & 0.431 & 0.444 & 0.506 & 0.585 & \textbf{0.593} & 0.593 \\
Depth 2 & 0.323 & 0.370 & 0.462 & 0.482 & 0.538 & 0.584 & 0.589 & \textbf{0.612} \\
\bottomrule
\end{tabular}%
}
\end{table*}

\begin{table}[t]
\centering
\caption{Hardware-ready Qiskit circuit path on a small MNIST-2 Aer fallback. This is not a hardware claim; it verifies the same circuit/count/readout pipeline used for IBM backend execution.}
\label{tab:hardware-ready}
\begin{tabular}{lcc}
\toprule
Method & Acc. & Bytes \\
\midrule
\method{}-Masked-HOP & 0.778 & 332 \\
CHOP & 0.756 & 268 \\
CProto & 0.753 & 140 \\
FedProto schema & 0.650 & 268 \\
FedAdam schema & 0.528 & 264 \\
No schema & 0.583 & 268 \\
Wrong schema & 0.561 & 396 \\
\bottomrule
\end{tabular}
\end{table}


The Qiskit circuit/count-to-readout path builds one circuit per observable and converts Aer bit counts into the same precomputed-readout format used by QProto. The robustness sweep reports five seeds per noise setting; the paired optimizer/backend comparisons use the three matched seeds available across all compared QFL baselines. Both are engineering compatibility checks, not hardware validation results.

\subsection{Additional Real Missing-View Fusion Check}

\begin{table}[t]
\centering
\caption{Strict equal-byte controls on the UCI HAR missing-view sensor benchmark. All rows use the same payload as CHOP.}
\label{tab:strict-equal-byte}
\begin{tabular}{lcc}
\toprule
Method & Acc. & Bytes \\
\midrule
Shared-view RFF equal bytes & 0.391 & 7708 \\
Zero-fill RFF equal bytes & 0.344 & 7708 \\
CProto equal bytes & 0.753 & 7708 \\
Group-balanced CProto equal bytes & 0.716 & 7708 \\
Random-key CProto equal bytes & 0.632 & 7708 \\
Random-key HOP equal bytes & 0.616 & 7708 \\
CHOP equal bytes & 0.632 & 7708 \\
Adaptive CHOP equal bytes & 0.752 & 7708 \\
\bottomrule
\end{tabular}
\end{table}

\begin{table}[t]
\centering
\caption{Strict equal-byte controls on the UCI Multiple Features multi-view benchmark. All rows use the same payload as CHOP.}
\label{tab:strict-equal-byte-mfeat}
\begin{tabular}{lcc}
\toprule
Method & Acc. & Bytes \\
\midrule
Shared-view RFF equal bytes & 0.694 & 16684 \\
Zero-fill RFF equal bytes & 0.766 & 16684 \\
CProto equal bytes & 0.868 & 16684 \\
Group-balanced CProto equal bytes & 0.876 & 16684 \\
Random-key CProto equal bytes & 0.891 & 16684 \\
Random-key HOP equal bytes & 0.883 & 16684 \\
CHOP equal bytes & 0.722 & 16684 \\
Adaptive CHOP equal bytes & 0.866 & 16684 \\
\bottomrule
\end{tabular}
\end{table}

\begin{table*}[t]
\centering
\caption{Real wearable multi-source missing-view fusion benchmark on UCI MHEALTH under a subject-held-out split. Windowed readout statistics are built from chest accelerometer, chest ECG, ankle IMU, and arm IMU sources. Equal-byte rows use the same per-client per-round payload as CHOP.}
\label{tab:if-mhealth-fusion}
\scriptsize
\setlength{\tabcolsep}{3pt}
\renewcommand{\arraystretch}{0.94}
\begin{tabular}{lcccc}
\toprule
Method & Acc. & 95\% CI & Bytes & Runtime (s) \\
\midrule
Mask-only metadata & 0.084 & 0.003 & 11092 & 0.580 \\
Zero-fill prototype & 0.789 & 0.024 & 11092 & 0.586 \\
Mean-impute prototype & 0.789 & 0.026 & 11092 & 0.705 \\
Value+mask prototype & 0.266 & 0.017 & 22132 & 1.229 \\
Anchor PCA imputation & 0.760 & 0.028 & 11092 & 11.689 \\
Learned ridge imputation & 0.828 & 0.028 & 11092 & 1.329 \\
Anchor autoencoder imputation & 0.804 & 0.029 & 11092 & 11.248 \\
Mask-aware pooled MLP & 0.860 & 0.020 & 22132 & 1.853 \\
HeMIS modality-dropout fusion & 0.889 & 0.025 & 22132 & 7.169 \\
View-wise late fusion & 0.814 & 0.026 & 11092 & 0.333 \\
Shared-view prototype & 0.719 & 0.044 & 1492 & 0.072 \\
Shared-view RFF equal bytes & 0.720 & 0.046 & 13876 & 0.864 \\
Zero-fill RFF equal bytes & 0.774 & 0.028 & 13876 & 1.005 \\
Coverage prototype all & 0.820 & 0.027 & 22132 & 1.406 \\
Group-balanced CProto all & 0.814 & 0.026 & 22132 & 1.644 \\
CProto equal bytes & 0.860 & 0.017 & 13876 & 0.917 \\
Group-balanced CProto equal bytes & 0.826 & 0.015 & 13876 & 1.034 \\
Random-key CProto equal bytes & 0.798 & 0.033 & 13876 & 0.625 \\
Random-key HOP equal bytes & 0.798 & 0.033 & 13876 & 1.209 \\
CHOP equal bytes & 0.865 & 0.041 & 13876 & 1.217 \\
Adaptive CHOP equal bytes & 0.878 & 0.022 & 13876 & 1.200 \\
\bottomrule
\end{tabular}
\end{table*}

\begin{table}[t]
\centering
\caption{Strict equal-byte controls on the real UCI MHEALTH wearable multi-source benchmark. All rows use the same payload as CHOP.}
\label{tab:strict-equal-byte-mhealth}
\begin{tabular}{lcc}
\toprule
Method & Acc. & Bytes \\
\midrule
Shared-view RFF equal bytes & 0.720 & 13876 \\
Zero-fill RFF equal bytes & 0.774 & 13876 \\
CProto equal bytes & 0.860 & 13876 \\
Group-balanced CProto equal bytes & 0.826 & 13876 \\
Random-key CProto equal bytes & 0.798 & 13876 \\
Random-key HOP equal bytes & 0.798 & 13876 \\
CHOP equal bytes & 0.865 & 13876 \\
Adaptive CHOP equal bytes & 0.878 & 13876 \\
\bottomrule
\end{tabular}
\end{table}

\input{tables/table_if_pamap2_fusion}
\input{tables/table_strict_equal_byte_pamap2}
\begin{table*}[t]
\centering
\caption{Real industrial multi-source missing-view fusion benchmark on UCI Hydraulic Systems condition monitoring (cooler target). Each sensor time series is summarized into cycle-level readout statistics, and each client observes a heterogeneous sensor subset. Equal-byte rows use the same per-client per-round payload as CHOP.}
\label{tab:if-hydraulic-fusion}
\scriptsize
\setlength{\tabcolsep}{3pt}
\renewcommand{\arraystretch}{0.94}
\begin{tabular}{lcccc}
\toprule
Method & Acc. & 95\% CI & Bytes & Runtime (s) \\
\midrule
Mask-only metadata & 0.333 & 0.000 & 2056 & 0.060 \\
Zero-fill prototype & 0.940 & 0.030 & 2056 & 0.056 \\
Mean-impute prototype & 0.939 & 0.031 & 2056 & 0.072 \\
Value+mask prototype & 0.667 & 0.095 & 4096 & 0.129 \\
Anchor PCA imputation & 0.980 & 0.006 & 2056 & 2.705 \\
Learned ridge imputation & 0.978 & 0.007 & 2056 & 0.533 \\
Anchor autoencoder imputation & 0.980 & 0.007 & 2056 & 5.384 \\
Mask-aware pooled MLP & 0.993 & 0.009 & 4096 & 0.420 \\
HeMIS modality-dropout fusion & 0.997 & 0.004 & 4096 & 5.429 \\
View-wise late fusion & 0.975 & 0.012 & 2056 & 0.015 \\
Shared-view prototype & 0.799 & 0.018 & 256 & 0.004 \\
Shared-view RFF equal bytes & 0.833 & 0.018 & 3472 & 0.159 \\
Zero-fill RFF equal bytes & 0.968 & 0.014 & 3472 & 0.259 \\
Coverage prototype all & 0.972 & 0.018 & 4096 & 0.130 \\
Group-balanced CProto all & 0.975 & 0.012 & 4096 & 0.153 \\
CProto equal bytes & 0.973 & 0.016 & 3472 & 0.107 \\
Group-balanced CProto equal bytes & 0.973 & 0.012 & 3472 & 0.116 \\
Random-key CProto equal bytes & 0.969 & 0.024 & 3472 & 0.068 \\
Random-key HOP equal bytes & 0.967 & 0.025 & 3472 & 0.133 \\
CHOP equal bytes & 0.977 & 0.016 & 3472 & 0.134 \\
\bottomrule
\end{tabular}
\end{table*}

\begin{table}[t]
\centering
\caption{Strict equal-byte controls on the real UCI Hydraulic Systems benchmark. All rows use the same payload as CHOP.}
\label{tab:strict-equal-byte-hydraulic}
\begin{tabular}{lcc}
\toprule
Method & Acc. & Bytes \\
\midrule
Shared-view RFF equal bytes & 0.833 & 3472 \\
Zero-fill RFF equal bytes & 0.968 & 3472 \\
CProto equal bytes & 0.973 & 3472 \\
Group-balanced CProto equal bytes & 0.973 & 3472 \\
Random-key CProto equal bytes & 0.969 & 3472 \\
Random-key HOP equal bytes & 0.967 & 3472 \\
CHOP equal bytes & 0.977 & 3472 \\
\bottomrule
\end{tabular}
\end{table}

\begin{table*}[t]
\centering
\caption{Real non-quantum missing-view fusion benchmark on the WDBC biomedical diagnosis dataset. Measurement families are treated as semantic views; each client observes a heterogeneous subset. Equal-byte rows use the same per-client per-round payload as CHOP.}
\label{tab:if-wdbc-fusion}
\scriptsize
\setlength{\tabcolsep}{3pt}
\renewcommand{\arraystretch}{0.94}
\begin{tabular}{lcccc}
\toprule
Method & Acc. & 95\% CI & Bytes & Runtime (s) \\
\midrule
Mask-only metadata & 0.509 & 0.015 & 252 & 0.003 \\
Zero-fill prototype & 0.921 & 0.006 & 252 & 0.003 \\
Mean-impute prototype & 0.922 & 0.006 & 252 & 0.004 \\
Value+mask prototype & 0.780 & 0.045 & 492 & 0.003 \\
Anchor PCA imputation & 0.882 & 0.014 & 252 & 0.565 \\
Learned ridge imputation & 0.924 & 0.008 & 252 & 0.029 \\
Anchor autoencoder imputation & 0.920 & 0.009 & 252 & 6.481 \\
Mask-aware pooled MLP & 0.920 & 0.008 & 492 & 0.706 \\
View-wise late fusion & 0.921 & 0.008 & 252 & 0.015 \\
Shared-view prototype & 0.896 & 0.009 & 36 & 0.004 \\
Shared-view RFF equal bytes & 0.896 & 0.006 & 396 & 0.010 \\
Zero-fill RFF equal bytes & 0.911 & 0.008 & 396 & 0.014 \\
Coverage prototype all & 0.922 & 0.006 & 492 & 0.017 \\
Group-balanced CProto all & 0.921 & 0.008 & 492 & 0.050 \\
CProto equal bytes & 0.922 & 0.006 & 396 & 0.011 \\
Group-balanced CProto equal bytes & 0.921 & 0.008 & 396 & 0.043 \\
Random-key HOP equal bytes & 0.905 & 0.011 & 396 & 0.017 \\
CHOP equal bytes & 0.923 & 0.006 & 396 & 0.019 \\
\bottomrule
\end{tabular}
\end{table*}

\begin{table}[t]
\centering
\caption{Strict equal-byte controls on the WDBC biomedical missing-view benchmark. All rows use the same payload as CHOP.}
\label{tab:strict-equal-byte-wdbc}
\begin{tabular}{lcc}
\toprule
Method & Acc. & Bytes \\
\midrule
Shared-view RFF equal bytes & 0.896 & 396 \\
Zero-fill RFF equal bytes & 0.911 & 396 \\
CProto equal bytes & 0.922 & 396 \\
Group-balanced CProto equal bytes & 0.921 & 396 \\
Random-key HOP equal bytes & 0.905 & 396 \\
CHOP equal bytes & 0.923 & 396 \\
\bottomrule
\end{tabular}
\end{table}

\begin{table}[t]
\centering
\caption{Strict equal-byte controls on the real Gas Sensor Array Drift benchmark. All rows use the same payload as CHOP.}
\label{tab:strict-equal-byte-gas-drift-random}
\begin{tabular}{lcc}
\toprule
Method & Acc. & Bytes \\
\midrule
Shared-view RFF equal bytes & 0.448 & 6940 \\
Zero-fill RFF equal bytes & 0.601 & 6940 \\
CProto equal bytes & 0.603 & 6940 \\
Group-balanced CProto equal bytes & 0.594 & 6940 \\
Random-key CProto equal bytes & 0.598 & 6940 \\
Random-key HOP equal bytes & 0.598 & 6940 \\
CHOP equal bytes & 0.592 & 6940 \\
Adaptive CHOP equal bytes & 0.604 & 6940 \\
\bottomrule
\end{tabular}
\end{table}

\begin{table}[t]
\centering
\caption{Strict equal-byte controls on the real Gas Sensor Array Drift benchmark. All rows use the same payload as CHOP.}
\label{tab:strict-equal-byte-gas-drift-temporal}
\begin{tabular}{lcc}
\toprule
Method & Acc. & Bytes \\
\midrule
Shared-view RFF equal bytes & 0.247 & 6940 \\
Zero-fill RFF equal bytes & 0.367 & 6940 \\
CProto equal bytes & 0.375 & 6940 \\
Group-balanced CProto equal bytes & 0.398 & 6940 \\
Random-key CProto equal bytes & 0.376 & 6940 \\
Random-key HOP equal bytes & 0.377 & 6940 \\
CHOP equal bytes & 0.364 & 6940 \\
Adaptive CHOP equal bytes & 0.364 & 6940 \\
\bottomrule
\end{tabular}
\end{table}


The first two tables isolate equal-byte rows from the UCI HAR and UCI Multiple Features main-text tables. The MHEALTH table gives the full wearable multi-source result summarized in the main text. PAMAP2 is included as a larger wearable boundary case where adaptive CHOP falls back to low-order CProto. Hydraulic Systems adds a real industrial condition-monitoring check: compact CProto/CHOP statistics are far above shared-view compression and close to PCA/ridge/late-fusion baselines, while HeMIS and mask-aware MLP remain stronger full-dimensional trainable references. WDBC is included as a biomedical breadth check. It is nearly saturated, so it is less diagnostic than UCI HAR, Multiple Features, MHEALTH, or Hydraulic Systems; the main text therefore treats it as a boundary case rather than a headline comparison. Gas Drift is reported as a negative boundary check: random and temporal splits both show that CHOP is not naturally favored there, even after cross-sensor HOP pair selection and a small HOP-weight screen.

\begin{figure}[t]
    \centering
    \includegraphics[width=\linewidth]{adaptive_chop_summary.pdf}
    \caption{Adaptive CHOP summary under matched communication. The mixture/fallback gate avoids pure-CHOP degradation on low-order tasks and gives the highest mean among compact equal-byte statistics on the ten-seed MHEALTH run.}
    \label{fig:adaptive-chop-supp}
\end{figure}

\begin{table*}[t]
\centering
\caption{Full paired significance tests for the main experimental claims. Reported values are mean paired accuracy deltas over matched seeds. One-sided paired \(t\)-tests test the directional claim; two-sided values are provided as a robustness check. Holm correction is applied across all reported comparisons.}
\label{tab:if-significance-full}
\scriptsize
\setlength{\tabcolsep}{2.5pt}
\renewcommand{\arraystretch}{0.78}
\resizebox{\textwidth}{!}{%
\begin{tabular}{llrrrrr}
\toprule
Setting & Comparison & \(n\) & \(\Delta\) acc. & \(d_z\) & Holm \(p_{1s}\) & Holm \(p_{2s}\) \\
\midrule
MNIST-10 & Masked-HOP - zero-fill & 5 & 0.052 & 4.67 & 0.0059 & 0.0119 \\
MNIST-10 & Masked-HOP - shared-observable & 5 & 0.050 & 4.78 & 0.0057 & 0.0113 \\
MNIST-10 & Masked-HOP - CProto & 5 & 0.007 & 2.98 & 0.0250 & 0.0500 \\
MNIST-10 & CHOP - zero-fill & 5 & 0.047 & 3.47 & 0.0149 & 0.0297 \\
Fashion-10 & Masked-HOP - zero-fill & 5 & 0.015 & 1.86 & 0.0851 & 0.1702 \\
Fashion-10 & Masked-HOP - shared-observable & 5 & 0.006 & 0.65 & 0.4422 & 0.8844 \\
Fashion-10 & Masked-HOP - CProto & 5 & 0.003 & 1.72 & 0.0869 & 0.1739 \\
Fashion-10 & CHOP - zero-fill & 5 & 0.013 & 1.53 & 0.1058 & 0.2116 \\
CIFAR-10 & Masked-HOP - zero-fill & 5 & 0.033 & 8.16 & 0.0009 & 0.0018 \\
CIFAR-10 & Masked-HOP - shared-observable & 5 & 0.029 & 1.83 & 0.0851 & 0.1702 \\
CIFAR-10 & Masked-HOP - CProto & 5 & 0.004 & 2.18 & 0.0574 & 0.1148 \\
CIFAR-10 & CHOP - zero-fill & 5 & 0.029 & 7.85 & 0.0010 & 0.0020 \\
Covariance sensing & CHOP - zero-fill & 5 & 0.105 & 3.54 & 0.0144 & 0.0289 \\
Covariance sensing & CHOP - shared-observable & 5 & 0.106 & 1.53 & 0.1058 & 0.2116 \\
Covariance sensing & CHOP - CProto & 5 & 0.374 & 10.35 & 0.0004 & 0.0007 \\
Covariance sensing & Masked-HOP - Masked & 5 & 0.319 & 8.40 & 0.0008 & 0.0016 \\
MNIST-10 & Masked-HOP - schema-FedProto & 5 & 0.332 & 11.56 & 0.0003 & 0.0005 \\
MNIST-10 & Masked-HOP - schema-FedAdam & 5 & 0.288 & 9.27 & 0.0006 & 0.0011 \\
MNIST-10 & Masked-HOP - no-schema & 5 & 0.525 & 19.85 & \(<0.0001\) & \(<0.0001\) \\
Fashion-10 & Masked-HOP - schema-FedProto & 5 & 0.140 & 6.17 & 0.0025 & 0.0050 \\
Fashion-10 & Masked-HOP - schema-FedAdam & 5 & 0.118 & 5.04 & 0.0049 & 0.0099 \\
Fashion-10 & Masked-HOP - no-schema & 5 & 0.423 & 23.67 & \(<0.0001\) & \(<0.0001\) \\
CIFAR-10 & Masked-HOP - schema-FedProto & 5 & 0.186 & 10.49 & 0.0004 & 0.0007 \\
CIFAR-10 & Masked-HOP - schema-FedAdam & 5 & 0.172 & 10.76 & 0.0003 & 0.0007 \\
CIFAR-10 & Masked-HOP - no-schema & 5 & 0.288 & 39.71 & \(<0.0001\) & \(<0.0001\) \\
PennyLane high-order & CHOP - CProto & 3 & 0.199 & 10.72 & 0.0255 & 0.0511 \\
PennyLane high-order & CHOP - schema-FedProto & 3 & 0.317 & 25.23 & 0.0063 & 0.0126 \\
Qiskit Aer & Masked-HOP - schema-FedProto & 3 & 0.169 & 2.91 & 0.1058 & 0.2116 \\
Qiskit Aer & Masked-HOP - FedDyn & 3 & 0.178 & 3.48 & 0.1058 & 0.2116 \\
UCI HAR & CProto eq. - zero-fill RFF eq. & 5 & 0.409 & 13.34 & 0.0002 & 0.0003 \\
UCI HAR & CProto eq. - shared-view RFF eq. & 5 & 0.362 & 24.84 & \(<0.0001\) & \(<0.0001\) \\
UCI HAR & Coverage all - anchor AE & 5 & 0.028 & 5.38 & 0.0040 & 0.0079 \\
MFeat & Group-CProto eq. - zero-fill RFF eq. & 10 & 0.110 & 1.89 & 0.0031 & 0.0062 \\
MFeat & Group-CProto eq. - shared-view RFF eq. & 10 & 0.182 & 4.17 & \(<0.0001\) & \(<0.0001\) \\
MFeat & Group-CProto all - mask-aware MLP & 10 & 0.008 & 1.58 & 0.0083 & 0.0166 \\
MHEALTH & CHOP eq. - CProto eq. & 10 & 0.005 & 0.09 & 1.0000 & 1.0000 \\
MHEALTH & CHOP eq. - zero-fill RFF eq. & 10 & 0.091 & 1.05 & 0.0579 & 0.1158 \\
MHEALTH & CHOP eq. - shared-view RFF eq. & 10 & 0.145 & 1.27 & 0.0255 & 0.0511 \\
MHEALTH & CHOP eq. - mask-aware MLP & 10 & 0.005 & 0.07 & 1.0000 & 1.0000 \\
MHEALTH & Adaptive CHOP - CProto eq. & 10 & 0.018 & 0.92 & 0.0869 & 0.1739 \\
MHEALTH & Random HOP - random CProto eq. & 10 & 0.000 & 0.09 & 1.0000 & 1.0000 \\
UCI HAR & Adaptive CHOP - CHOP eq. & 5 & 0.120 & 4.82 & 0.0057 & 0.0113 \\
MFeat & Adaptive CHOP - CHOP eq. & 10 & 0.144 & 1.62 & 0.0071 & 0.0142 \\
PAMAP2 & Adaptive CHOP - CHOP eq. & 5 & 0.230 & 2.74 & 0.0269 & 0.0537 \\
PAMAP2 & Adaptive CHOP - zero-fill RFF eq. & 5 & 0.116 & 2.92 & 0.0255 & 0.0511 \\
PAMAP2 & Adaptive CHOP - shared-view RFF eq. & 5 & 0.450 & 16.60 & \(<0.0001\) & 0.0001 \\
\bottomrule
\end{tabular}%
}
\end{table*}


\paragraph{Stability and claim boundaries.}
The seed-level tables should be read with the scoped claims in Table~\ref{tab:claim-evidence-map}. The most stable comparisons are the schema-correct rows against no-schema, wrong-schema, shared-view, and equal-byte RFF controls. UCI HAR, MHEALTH, and Multiple Features also show stable advantages over equal-byte shared-view compression. Some comparisons are intentionally not used as headline claims: MFeat random-key CProto is strong, so that benchmark supports schema-aware low-order fusion rather than selected-key superiority; adaptive CHOP on MHEALTH has the highest mean among compact equal-byte statistics but is reported as interaction-heavy evidence rather than as a universal HOP guarantee; and Qiskit Aer rows are compatibility checks with limited matched seeds, not hardware validation.

\subsection{High-Order and Key-Policy Diagnostics}

\begin{table*}[t]
\centering
\caption{High-order readout experiments. The PennyLane benchmark uses circuit-generated readouts with class means equalized, so the dominant signal is second-order.}
\label{tab:masked-hop-highorder}
\resizebox{\textwidth}{!}{%
\begin{tabular}{lcccc}
\toprule
Setting & Method & Acc. & 95\% CI & Bytes \\
\midrule
PennyLane high-order & \method{}-Masked-HOP & 0.550 & 0.009 & 3476 \\
PennyLane high-order & CHOP, \(K=24\) & 0.602 & 0.061 & 2324 \\
PennyLane high-order & \method{}-Masked & 0.422 & 0.044 & 1940 \\
PennyLane high-order & CProto, \(K=24\) & 0.403 & 0.061 & 788 \\
PennyLane high-order & shared-observable & 0.446 & 0.053 & 1556 \\
PennyLane high-order & old \hopmethod & 0.320 & 0.028 & 3092 \\
PennyLane high-order & no-schema & 0.338 & 0.057 & 1556 \\
\midrule
Synthetic covariance & \method{}-Masked-HOP & 0.808 & 0.083 & 4116 \\
Synthetic covariance & \method{}-Masked & 0.467 & 0.085 & 2068 \\
Synthetic covariance & old \hopmethod & 0.795 & 0.059 & 2260 \\
Synthetic covariance & low-order prototype & 0.449 & 0.115 & 212 \\
Synthetic covariance & wrong-schema HOP & 0.445 & 0.029 & 2260 \\
\bottomrule
\end{tabular}%
}
\end{table*}

\begin{table}[t]
\centering
\caption{CHOP communication sweeps. CProto uses the same selected observable keys but omits HOP.}
\label{tab:chop-sweep}
\begin{tabular}{llccc}
\toprule
Setting & Variant & \(K\) & Acc. & Bytes \\
\midrule
PennyLane high-order & CProto & 8 & 0.389 & 276 \\
PennyLane high-order & CHOP & 8 & 0.639 & 1812 \\
PennyLane high-order & CHOP & 24 & 0.602 & 2324 \\
PennyLane high-order & CHOP & 48 & 0.579 & 3092 \\
\midrule
MNIST-10 & CHOP & 32 & 0.663 & 6444 \\
MNIST-10 & CHOP & 64 & 0.730 & 9004 \\
MNIST-10 & CHOP & 96 & 0.751 & 11564 \\
MNIST-10 & CHOP & 128 & 0.753 & 14124 \\
\midrule
Qiskit Aer MNIST-4 & CHOP & 4 & 0.556 & 404 \\
Qiskit Aer MNIST-4 & CHOP & 8 & 0.579 & 532 \\
Qiskit Aer MNIST-4 & CHOP & 10 & 0.594 & 596 \\
Qiskit Aer MNIST-4 & CHOP & 12 & 0.597 & 660 \\
\bottomrule
\end{tabular}
\end{table}

\begin{table*}[t]
\centering
\caption{CHOP key-selection controls. Anchor-informed key policies are compared to random keys at the same \(K\).}
\label{tab:chop-key-policy}
\resizebox{\textwidth}{!}{%
\begin{tabular}{llccc}
\toprule
Setting & Key policy & Method & Acc. & Bytes \\
\midrule
PennyLane high-order, \(K=8\) & coverage & CHOP & 0.659 & 1812 \\
PennyLane high-order, \(K=8\) & variance & CHOP & 0.624 & 1812 \\
PennyLane high-order, \(K=8\) & random & CHOP & 0.431 & 1812 \\
PennyLane high-order, \(K=8\) & coverage & CProto & 0.394 & 276 \\
\midrule
MNIST-10, \(K=96\) & variance & CHOP & 0.751 & 11564 \\
MNIST-10, \(K=96\) & coverage & CHOP & 0.749 & 11564 \\
MNIST-10, \(K=96\) & random & CHOP & 0.662 & 11564 \\
\midrule
CIFAR-10, \(K=96\) & coverage & CHOP & 0.452 & 11564 \\
CIFAR-10, \(K=96\) & variance & CHOP & 0.451 & 11564 \\
CIFAR-10, \(K=96\) & random & CHOP & 0.391 & 11564 \\
\midrule
Qiskit Aer, \(K=10\) & variance & CHOP & 0.594 & 596 \\
Qiskit Aer, \(K=10\) & coverage & CHOP & 0.590 & 596 \\
Qiskit Aer, \(K=10\) & random & CHOP & 0.582 & 596 \\
\bottomrule
\end{tabular}%
}
\end{table*}

\begin{table*}[t]
\centering
\caption{Paired seed-level accuracy deltas for the main high-order and key-selection claims.}
\label{tab:pairwise-deltas}
\resizebox{\textwidth}{!}{%
\begin{tabular}{lccc}
\toprule
Comparison & \(n\) & Mean delta & Wins \\
\midrule
Masked-HOP minus Masked, PennyLane high-order & 3 & 0.128 \(\pm\) 0.042 & 3/3 \\
Masked-HOP minus old HOP, PennyLane high-order & 3 & 0.230 \(\pm\) 0.035 & 3/3 \\
CHOP minus CProto, PennyLane high-order \(K=8\) & 3 & 0.250 \(\pm\) 0.071 & 3/3 \\
Coverage-key CHOP minus random-key CHOP, \(K=8\) & 5 & 0.228 \(\pm\) 0.066 & 5/5 \\
Variance-key CHOP minus random-key CHOP, MNIST-10 & 3 & 0.089 \(\pm\) 0.016 & 3/3 \\
Coverage-key CHOP minus random-key CHOP, CIFAR-10 & 3 & 0.061 \(\pm\) 0.044 & 3/3 \\
CHOP minus schema-FedProto, MNIST-10 & 3 & 0.333 \(\pm\) 0.027 & 3/3 \\
CHOP minus schema-FedProto, CIFAR-10 & 3 & 0.179 \(\pm\) 0.058 & 3/3 \\
CHOP minus schema-FedProto, Qiskit Aer & 3 & 0.151 \(\pm\) 0.227 & 3/3 \\
CHOP minus schema-FedProto, PennyLane high-order & 3 & 0.317 \(\pm\) 0.031 & 3/3 \\
CHOP minus schema-FedAdam, MNIST-10 & 3 & 0.291 \(\pm\) 0.031 & 3/3 \\
CHOP minus schema-FedAdam, CIFAR-10 & 3 & 0.169 \(\pm\) 0.054 & 3/3 \\
CHOP minus schema-FedAdam, PennyLane high-order & 3 & 0.286 \(\pm\) 0.042 & 3/3 \\
\bottomrule
\end{tabular}%
}
\end{table*}


\begin{figure}[t]
    \centering
    \includegraphics[width=\linewidth]{chop_sweep_accuracy.pdf}
    \caption{CHOP communication sweeps. CHOP and CProto use the same selected observable keys; the gap isolates mask-aware high-order statistics.}
    \label{fig:chop-sweep}
\end{figure}

\subsection{Schema, Shot, and Communication Diagnostics}

\begin{table}[t]
\centering
\caption{Schema comparability benchmark on synthetic heterogeneous readout.}
\label{tab:schema}
\begin{tabular}{lccc}
\toprule
Method & Acc. & Std. & Bal. acc. \\
\midrule
Full + shrinkage & 0.621 & 0.055 & 0.620 \\
No schema & 0.422 & 0.047 & 0.421 \\
Forced canonical & 0.415 & 0.019 & 0.415 \\
Wrong schema & 0.290 & 0.010 & 0.292 \\
\bottomrule
\end{tabular}
\end{table}

\begin{table}[t]
\centering
\caption{High-order outer-product benchmark on covariance-structured data.}
\label{tab:hop}
\begin{tabular}{lccc}
\toprule
Method & Acc. & Std. & Bal. acc. \\
\midrule
Full + shrinkage & 0.795 & 0.027 & 0.799 \\
QProto-HOP & 0.795 & 0.024 & 0.798 \\
Prototype & 0.449 & 0.046 & 0.449 \\
Wrong schema & 0.445 & 0.012 & 0.445 \\
No schema & 0.308 & 0.013 & 0.307 \\
\bottomrule
\end{tabular}
\end{table}

\begin{table}[t]
\centering
\caption{CIFAR-4 shot-noise sweep. Results are means over 5 seeds.}
\label{tab:shot}
\begin{tabular}{lccccc}
\toprule
Shots & Full & Proto & No schema & Wrong & FedAvg \\
\midrule
32 & 0.474 & 0.469 & 0.334 & 0.261 & 0.286 \\
64 & 0.476 & 0.471 & 0.336 & 0.261 & 0.286 \\
128 & 0.475 & 0.472 & 0.337 & 0.262 & 0.287 \\
256 & 0.474 & 0.472 & 0.337 & 0.262 & 0.287 \\
1024 & 0.473 & 0.472 & 0.338 & 0.262 & 0.287 \\
\bottomrule
\end{tabular}
\end{table}

\begin{table}[t]
\centering
\caption{Matched communication on the high-order synthetic benchmark.}
\label{tab:comm-synth}
\begin{tabular}{lccc}
\toprule
Budget & Method & Bytes & Acc. \\
\midrule
Low & Prototype & 212 & 0.449 \\
Low & HOP & 212 & 0.504 \\
High & Prototype & 2260 & 0.749 \\
High & HOP & 2260 & 0.795 \\
\bottomrule
\end{tabular}
\end{table}

\begin{table}[t]
\centering
\caption{CIFAR-4 matched communication curve.}
\label{tab:comm-real}
\begin{tabular}{lcc}
\toprule
Bytes & Prototype & HOP \\
\midrule
1556 & 0.446 & 0.429 \\
3092 & 0.488 & 0.446 \\
4116 & 0.496 & 0.468 \\
\bottomrule
\end{tabular}
\end{table}


\begin{figure}[t]
    \centering
    \includegraphics[width=\linewidth]{schema_triptych_accuracy.pdf}
    \caption{Schema comparability benchmark. The gain comes from making readout coordinates semantically comparable before aggregation.}
    \label{fig:schema}
\end{figure}

\begin{figure}[t]
    \centering
    \includegraphics[width=\linewidth]{hop_ablation_accuracy.pdf}
    \caption{High-order benchmark. HOP separates covariance-structured classes that low-order prototypes collapse.}
    \label{fig:hop}
\end{figure}

\begin{figure}[t]
    \centering
    \includegraphics[width=\linewidth]{real_cifar4_shot_sweep.pdf}
    \caption{CIFAR-4 shot sweep under heterogeneous readout.}
    \label{fig:shot}
\end{figure}

\begin{figure}[t]
    \centering
    \includegraphics[width=\linewidth]{real_cifar4_comm_budget.pdf}
    \caption{CIFAR-4 communication curve. On this low-order-separable real feature task, HOP is not the primary source of gain.}
    \label{fig:comm-real}
\end{figure}


\section{Additional Theory Details}

\subsection{Plug-In Classification Consistency}

Let \(G(z)=g(A_{\schema}(r(x)))\) denote the protocol feature used by the classifier, where \(g\) may contain low-order RFF features, coverage-aware values, and HOP features. Define the population class statistic
\begin{equation}
    \theta_c^\star=\E[G(z)\mid y=c]
\end{equation}
and the population score
\begin{equation}
    s_c^\star(z)=\|G(z)-\theta_c^\star\|_2^2.
\end{equation}
The empirical server statistic \(\hat{\theta}_c\) is a plug-in estimator of \(\theta_c^\star\). If the margin
\begin{equation}
    \Delta(z)=\min_{c\neq y}\{s_c^\star(z)-s_y^\star(z)\}
\end{equation}
is positive and
\begin{equation}
    \max_c | \hat{s}_c(z)-s_c^\star(z)| < \Delta(z)/2,
\end{equation}
then the empirical classifier returns the same label as the population nearest-prototype classifier. This is the margin condition used implicitly in the finite-budget theorem in the main paper. It converts estimation error in comparable readout statistics into a classification guarantee.

\subsection{Finite-Shot Constants}

For a Pauli observable with outcomes in \(\{-1,+1\}\), Hoeffding's inequality gives
\begin{equation}
    \Pr(|\hat{\mu}_{ia}-\mu_{ia}|>\epsilon)
    \le 2\exp(-S_i\epsilon^2/2).
\end{equation}
After a union bound over \(D\) public observables, with probability at least \(1-\delta\),
\begin{equation}
    \|\hat{v}-v\|_2
    \le
    \sqrt{\frac{2D\log(2D/\delta)}{S_{\min}}}.
\end{equation}
If the feature map is \(L_g\)-Lipschitz and bounded by \(R_g\), the induced score perturbation is bounded by
\begin{equation}
    |\|\hat{G}-\theta_c\|_2^2-\|G-\theta_c\|_2^2|
    \le
    2R_gL_g\|\hat{v}-v\|_2+L_g^2\|\hat{v}-v\|_2^2.
\end{equation}
This explains why shot noise enters the finite-budget theorem through a \(D^{1/2}S_{\min}^{-1/2}\) term, while schema mismatch does not vanish with shots.

\subsection{Coverage-Aware Estimator Details}

The coverage-aware estimator stores two sufficient statistics per class and source coordinate: a value sum and an observation count. This is equivalent to estimating
\begin{equation}
    \E[V_j\mid y=c,M_j=1]
\end{equation}
instead of estimating \(\E[M_jV_j\mid y=c]\). Zero-fill prototypes estimate the latter quantity, which mixes measurement value and observation probability. This distinction is important when missingness is caused by source availability or measurement design: a missing observable should not be interpreted as a measured zero expectation. If missingness is conditionally independent of the value given the class, the coverage-aware estimator targets the full class-conditional coordinate mean. If missingness is informative, it targets the observed-view mean and the mask/coverage pattern remains available to the classifier.

\subsection{Schema-Uncertainty Details}

The main schema mismatch theorem treats a fully wrong permutation for clarity. A partially wrong schema produces the same kind of non-vanishing semantic bias. If a fraction \(\rho\) of coordinate mass is assigned to incorrect public keys, boundedness gives an \(O(\rho)\) bias term in the affected class prototypes. This term is separate from sample, shot, and random-feature errors. It can only be reduced by improving schema metadata, using calibration anchors to detect mismatches, or modeling schema uncertainty explicitly. The current method assumes protocol-defined observable identities are available; noisy schema inference is left for future work.

\subsection{Random-Feature and Communication Terms}

For bounded shift-invariant kernels, random Fourier features estimate kernel values with Monte Carlo error \(O(d^{-1/2})\) for fixed pairs. A uniform statement over \(M\) evaluated pairs follows by a union bound:
\begin{equation}
    \max_{\ell\le M}|\hat{k}(z_\ell,z'_\ell)-k(z_\ell,z'_\ell)|
    =
    O\!\left(\sqrt{\frac{\log(M/\delta)}{d}}\right).
\end{equation}
Because each class prototype uses \(d\) low-order coordinates and \(h\) high-order coordinates, the payload scales as \(C(d+h)\). Thus, increasing communication reduces sketch error but does not reduce schema error. This separation is central to the protocol claim: no amount of RFF or HOP dimension can fix an incorrect observable map.

\subsection{HOP Advantage Under Equal Means}

Suppose two classes have equal adapted readout means but covariance difference \(A=\Sigma_1-\Sigma_2\neq0\). For an isotropic Gaussian projection \(p\),
\begin{equation}
    \E_p[(p^\top A p)^2]
    =2\|A\|_F^2+(\operatorname{tr} A)^2.
\end{equation}
Therefore, whenever the class covariance matrices differ, projected quadratic features have nonzero expected squared gap even when low-order mean prototypes coincide. With \(h\) independent HOP projections, the probability that every projected gap is below a fixed detectable threshold decays exponentially in \(h\) under standard small-ball assumptions. This is the formal reason CHOP is strong on the covariance-sensing and PennyLane high-order tasks but only mildly helpful on low-order image tasks.

\subsection{Group-Balanced Source Scoring}

The group-balanced CProto rule is used when public coordinates are partitioned into semantic source groups \(G_1,\ldots,G_Q\), such as sensor blocks or feature families. Let
\begin{equation}
    \ell_{cq}(x)=
    \frac{\sum_{j\in G_q}m_j(x)\mathbf{1}\{N_{cj}>0\}w_{cj}(v_j(x)-\bar v_{cj})^2}
    {\sum_{j\in G_q}m_j(x)\mathbf{1}\{N_{cj}>0\}w_{cj}+\epsilon}
\end{equation}
be the within-group coverage-normalized class distance. Group-balanced scoring uses
\begin{equation}
    d_c^{\mathrm{grp}}(x)=|\mathcal Q_x|^{-1}\sum_{q\in\mathcal Q_x}\ell_{cq}(x),
\end{equation}
where \(\mathcal Q_x\) is the set of observed semantic groups.

\begin{proposition}[Dimensional dominance of coordinate-wise scoring]
Assume each semantic group \(G_q\) has an expected within-group loss \(L_{cq}=\E[\ell_{cq}(X)]\), and suppose a coordinate-wise distance averages over all observed coordinates without group normalization. If all coordinates in a group have comparable expected loss, the coordinate-wise score estimates a dimension-weighted objective proportional to
\begin{equation}
    \sum_q |G_q|L_{cq},
\end{equation}
whereas group-balanced scoring estimates a source-balanced objective proportional to
\begin{equation}
    \sum_q L_{cq}.
\end{equation}
Thus, when source groups have very different dimensionalities, coordinate-wise scoring can be dominated by the largest source even if the smaller sources are semantically equally important.
\end{proposition}

\begin{proof}
Ignoring the small \(\epsilon\) term and conditioning on the observed groups, coordinate-wise scoring is an average of per-coordinate squared losses over \(\cup_qG_q\). If the average coordinate loss within group \(q\) is \(L_{cq}\), then the numerator contributes approximately \(|G_q|L_{cq}\). Summing over groups gives a dimension-weighted objective. Group-balanced scoring first averages within \(G_q\) and then averages over groups, so each observed semantic source contributes one group-level term. This is the intended objective in multi-source fusion settings where a source's importance should not be determined only by how many exported coordinates it has.
\end{proof}

\subsection{Adaptive CHOP Calibration Bound}

Adaptive CHOP chooses among a finite candidate set \(\mathcal A\), containing the low-order CProto branch and a small grid of CHOP mixtures. The choice is made on a training-only calibration split and the selected candidate is then refit on the full training set. The following bound clarifies what this gate can and cannot guarantee.

\begin{proposition}[Finite calibration selection]
Let \(\mathcal A\) contain \(M\) candidate scoring rules with true calibration accuracies \(a(\alpha)\) and empirical calibration accuracies \(\hat a(\alpha)\) computed on \(n_{\mathrm{cal}}\) independent calibration examples. Let \(\hat\alpha=\arg\max_{\alpha\in\mathcal A}\hat a(\alpha)\). With probability at least \(1-\delta\),
\begin{equation}
    a(\hat\alpha)
    \ge
    \max_{\alpha\in\mathcal A}a(\alpha)
    -2\sqrt{\frac{\log(2M/\delta)}{2n_{\mathrm{cal}}}} .
\end{equation}
If the protocol requires a margin \(\tau\) over the low-order branch before enabling CHOP, the selected high-order branch has the same guarantee relative to the best candidate in \(\mathcal A\), while false high-order activation is controlled by the same uniform deviation term plus \(\tau\).
\end{proposition}

\begin{proof}
For a fixed candidate, Hoeffding's inequality gives
\begin{equation}
    \Pr(|\hat a(\alpha)-a(\alpha)|>\epsilon)\le2\exp(-2n_{\mathrm{cal}}\epsilon^2).
\end{equation}
A union bound over \(M\) candidates gives uniform deviation at level
\(\epsilon=\sqrt{\log(2M/\delta)/(2n_{\mathrm{cal}})}\). Let \(\alpha^\star=\arg\max_{\alpha}a(\alpha)\). Then
\begin{equation}
    a(\hat\alpha)\ge \hat a(\hat\alpha)-\epsilon
    \ge \hat a(\alpha^\star)-\epsilon
    \ge a(\alpha^\star)-2\epsilon.
\end{equation}
The margin rule compares the best CHOP mixture to the low-order candidate using the same uniformly controlled empirical accuracies, so the probability of enabling a worse high-order branch is bounded by the same calibration uncertainty unless the true gap lies within the margin band. This is why the method treats adaptive CHOP as a train-calibrated fallback rule, not as a universal high-order guarantee.
\end{proof}

\subsection{Combined Finite-Budget Error Decomposition}

The main theorems separate semantic error, sampling error, finite-shot error, sketch approximation, and high-order allocation. For a compact way to read them together, write \(\Pi\) for the schema used by the server and let \(\Pi^\star\) be the true semantic map. For a fixed class and query, the score error can be decomposed as
\begin{equation}
\begin{aligned}
|\hat s_c(z)-s_c^\star(z)|
\le&
\underbrace{E_{\mathrm{schema}}(\Pi,\Pi^\star)}_{\text{semantic bias}}
+\underbrace{E_{\mathrm{sample}}(n_c)}_{\text{finite examples}}
+\underbrace{E_{\mathrm{shot}}(S_{\min})}_{\text{finite readout}}\\
&+\underbrace{E_{\mathrm{rff}}(d)}_{\text{low-order sketch}}
+\underbrace{E_{\mathrm{hop}}(h)}_{\text{high-order sketch}}
+\underbrace{E_{\mathrm{cal}}(M,n_{\mathrm{cal}})}_{\text{adaptive gate}} .
\end{aligned}
\end{equation}
Under the assumptions in the main text, these terms may be chosen with high probability as
\begin{equation}
\begin{aligned}
    E_{\mathrm{sample}}&=O(n_c^{-1/2}),&
    E_{\mathrm{shot}}&=O(D^{1/2}S_{\min}^{-1/2}),\\
    E_{\mathrm{rff}}&=O(d^{-1/2}),&
    E_{\mathrm{hop}}&=O(h^{-1/2}),
\end{aligned}
\end{equation}
and
\begin{equation}
    E_{\mathrm{cal}}=
    O\!\left(\sqrt{\frac{\log M+\log(1/\delta)}{n_{\mathrm{cal}}}}\right).
\end{equation}
The important qualitative distinction is that \(E_{\mathrm{schema}}\) does not vanish with more samples, shots, or sketch dimensions unless the schema itself is corrected. This is the mathematical version of the paper's central information-fusion claim: communication and estimation can refine comparable evidence, but they cannot repair anonymous or wrongly named coordinates.
